% Options for packages loaded elsewhere
\PassOptionsToPackage{unicode}{hyperref}
\PassOptionsToPackage{hyphens}{url}
\documentclass[
]{article}
\usepackage{xcolor}
\usepackage{amsmath,amssymb}
\setcounter{secnumdepth}{-\maxdimen} % remove section numbering
\usepackage{iftex}
\ifPDFTeX
  \usepackage[T1]{fontenc}
  \usepackage[utf8]{inputenc}
  \usepackage{textcomp} % provide euro and other symbols
\else % if luatex or xetex
  \usepackage{unicode-math} % this also loads fontspec
  \defaultfontfeatures{Scale=MatchLowercase}
  \defaultfontfeatures[\rmfamily]{Ligatures=TeX,Scale=1}
\fi
\usepackage{lmodern}
\ifPDFTeX\else
  % xetex/luatex font selection
\fi
% Use upquote if available, for straight quotes in verbatim environments
\IfFileExists{upquote.sty}{\usepackage{upquote}}{}
\IfFileExists{microtype.sty}{% use microtype if available
  \usepackage[]{microtype}
  \UseMicrotypeSet[protrusion]{basicmath} % disable protrusion for tt fonts
}{}
\makeatletter
\@ifundefined{KOMAClassName}{% if non-KOMA class
  \IfFileExists{parskip.sty}{%
    \usepackage{parskip}
  }{% else
    \setlength{\parindent}{0pt}
    \setlength{\parskip}{6pt plus 2pt minus 1pt}}
}{% if KOMA class
  \KOMAoptions{parskip=half}}
\makeatother
\usepackage{color}
\usepackage{fancyvrb}
\newcommand{\VerbBar}{|}
\newcommand{\VERB}{\Verb[commandchars=\\\{\}]}
\DefineVerbatimEnvironment{Highlighting}{Verbatim}{commandchars=\\\{\}}
% Add ',fontsize=\small' for more characters per line
\newenvironment{Shaded}{}{}
\newcommand{\AlertTok}[1]{\textcolor[rgb]{1.00,0.00,0.00}{\textbf{#1}}}
\newcommand{\AnnotationTok}[1]{\textcolor[rgb]{0.38,0.63,0.69}{\textbf{\textit{#1}}}}
\newcommand{\AttributeTok}[1]{\textcolor[rgb]{0.49,0.56,0.16}{#1}}
\newcommand{\BaseNTok}[1]{\textcolor[rgb]{0.25,0.63,0.44}{#1}}
\newcommand{\BuiltInTok}[1]{\textcolor[rgb]{0.00,0.50,0.00}{#1}}
\newcommand{\CharTok}[1]{\textcolor[rgb]{0.25,0.44,0.63}{#1}}
\newcommand{\CommentTok}[1]{\textcolor[rgb]{0.38,0.63,0.69}{\textit{#1}}}
\newcommand{\CommentVarTok}[1]{\textcolor[rgb]{0.38,0.63,0.69}{\textbf{\textit{#1}}}}
\newcommand{\ConstantTok}[1]{\textcolor[rgb]{0.53,0.00,0.00}{#1}}
\newcommand{\ControlFlowTok}[1]{\textcolor[rgb]{0.00,0.44,0.13}{\textbf{#1}}}
\newcommand{\DataTypeTok}[1]{\textcolor[rgb]{0.56,0.13,0.00}{#1}}
\newcommand{\DecValTok}[1]{\textcolor[rgb]{0.25,0.63,0.44}{#1}}
\newcommand{\DocumentationTok}[1]{\textcolor[rgb]{0.73,0.13,0.13}{\textit{#1}}}
\newcommand{\ErrorTok}[1]{\textcolor[rgb]{1.00,0.00,0.00}{\textbf{#1}}}
\newcommand{\ExtensionTok}[1]{#1}
\newcommand{\FloatTok}[1]{\textcolor[rgb]{0.25,0.63,0.44}{#1}}
\newcommand{\FunctionTok}[1]{\textcolor[rgb]{0.02,0.16,0.49}{#1}}
\newcommand{\ImportTok}[1]{\textcolor[rgb]{0.00,0.50,0.00}{\textbf{#1}}}
\newcommand{\InformationTok}[1]{\textcolor[rgb]{0.38,0.63,0.69}{\textbf{\textit{#1}}}}
\newcommand{\KeywordTok}[1]{\textcolor[rgb]{0.00,0.44,0.13}{\textbf{#1}}}
\newcommand{\NormalTok}[1]{#1}
\newcommand{\OperatorTok}[1]{\textcolor[rgb]{0.40,0.40,0.40}{#1}}
\newcommand{\OtherTok}[1]{\textcolor[rgb]{0.00,0.44,0.13}{#1}}
\newcommand{\PreprocessorTok}[1]{\textcolor[rgb]{0.74,0.48,0.00}{#1}}
\newcommand{\RegionMarkerTok}[1]{#1}
\newcommand{\SpecialCharTok}[1]{\textcolor[rgb]{0.25,0.44,0.63}{#1}}
\newcommand{\SpecialStringTok}[1]{\textcolor[rgb]{0.73,0.40,0.53}{#1}}
\newcommand{\StringTok}[1]{\textcolor[rgb]{0.25,0.44,0.63}{#1}}
\newcommand{\VariableTok}[1]{\textcolor[rgb]{0.10,0.09,0.49}{#1}}
\newcommand{\VerbatimStringTok}[1]{\textcolor[rgb]{0.25,0.44,0.63}{#1}}
\newcommand{\WarningTok}[1]{\textcolor[rgb]{0.38,0.63,0.69}{\textbf{\textit{#1}}}}
\setlength{\emergencystretch}{3em} % prevent overfull lines
\providecommand{\tightlist}{%
  \setlength{\itemsep}{0pt}\setlength{\parskip}{0pt}}
\usepackage{bookmark}
\IfFileExists{xurl.sty}{\usepackage{xurl}}{} % add URL line breaks if available
\urlstyle{same}
\hypersetup{
  pdftitle={U-AUB-02: bidefect count route{,} current formal frontier},
  pdfauthor={Kieran McShane},
  hidelinks,
  pdfcreator={LaTeX via pandoc}}

\title{U-AUB-02: bidefect count route, current formal frontier}
\author{Kieran McShane}
\date{2026-06-07}

\begin{document}
\maketitle

\section{U-AUB-02: bidefect count route, current formal
frontier}\label{u-aub-02-bidefect-count-route-current-formal-frontier}

\textbf{Status.} U-AUB-02 is still open. The current Lean declaration is
a checked conditional reduction: it proves that the listed Biane,
frontier, and Chapuy inputs imply \texttt{AubrunBiDefectCountBound\ m};
it does not prove those inputs.

\subsection{Target}\label{target}

The formal target behind U-AUB-02 is the corrected Aubrun count input
used by the expectation pipeline. In mathematical terms, the route now
tries to prove a bidefect count theorem of the following form.

For each Wick length \(m\), every admissible pair of defects \((a,b)\)
should satisfy \[
  \#\{\pi : \delta_+(\pi)=a,\ \delta_-(\pi)=b\}
  \le
  E_m(a,b),
\] where \(E_m(a,b)\) is the explicit Aubrun envelope used later in the
scalar sum. Once this count is known, the already checked scalar
absorption turns the double defect sum into shifted bases rather than a
false common-base maximum.

The remaining theorem-strength work is exactly the following four
groups:

\begin{enumerate}
\def\labelenumi{\arabic{enumi}.}
\tightlist
\item
  the large-gap left Biane adjacent-link theorem;
\item
  the large-gap right Biane adjacent-link theorem;
\item
  the side-specific residual branch collision theorems for non-fixed
  predecessor skips, including the nonwrapping smaller
  internal-subinterval child branch;
\item
  the Chapuy/trisection recurrence.
\end{enumerate}

The active Lean endpoint is the checked reduction from those four groups
to the bidefect count bound:

\begin{Shaded}
\begin{Highlighting}[]
\NormalTok{AppendixB.AubrunBiDefectCountBound\_of\_bianeAdjacentLinkHalves\_left\_gap\_three\_start\_impossible\_frontiers\_nestedStrictReturnIntervalOpenNonwrapOffCycleSmallerInternalOnly\_and\_chapuy}
\end{Highlighting}
\end{Shaded}

This name should not be read as a solved ticket. Its theorem statement
still contains the theorem-strength hypotheses listed immediately below.

The previous reduction layer was:

\begin{Shaded}
\begin{Highlighting}[]
\NormalTok{AppendixB.AubrunBiDefectCountBound\_of\_bianeAdjacentLinkHalves\_left\_gap\_three\_start\_impossible\_frontiers\_skipMinimalFrontierCommonCycle\_and\_chapuy}
\end{Highlighting}
\end{Shaded}

Delta: the previous layers still asked the collision theorem to work
from the minimal-frontier packet. The active layer uses
\texttt{permCyclicPredecessorSkip\_direct\_commonCycle} internally,
supplies branch-normalized return intervals internally, descends
internal-right repeats to actual shorter subintervals, splits the three
origin-tag side cases, and exposes the nonwrapping internal-subinterval
branch only in the off-cycle open-internal form. The newest checked
support then sends the off-cycle witness through a smaller side-aware
residual child and splits that child into primitive and
internal-subinterval branches, then absorbs the strict primitive child
alternatives into the already theorem-facing parent primitive residual
leaf.

Recent left-slice improvement: the residual branch
\texttt{not\ pi(a,a+2)\ and\ S(a,a+2)} is contradictory. The repeat
\texttt{S(a,a+2)} copies the one-step contour move from \texttt{a+2} to
\texttt{c} into a length-two path from \texttt{S(a)} to \texttt{S(c)}.
Tree uniqueness compares this copied path with the direct length-two
path from \texttt{pi(a,c)} and forces \texttt{T(a+2)=T(a)}, i.e.
\texttt{pi(a,a+2)}, contradicting the branch assumption.

That change is \texttt{U-AUB-316} in the ledger. It removes the last
left-start gap-three parameter; the start slice \texttt{c-a=3,\ b=a+1}
is no longer present on the endpoint surface. The later
\texttt{U-AUB-317} closes the right special-slice branch
\texttt{S(b,c)}: if \texttt{S(c,d)} failed, copying the contour step
\texttt{c\ -\textgreater{}\ d} across \texttt{S(b)=S(c)} would give a
simple path from \texttt{S(b)} to \texttt{S(d)}, and tree uniqueness
against the direct path from \texttt{pi(b,d)} would force
\texttt{pi(b,c)}, contradicting \texttt{pi(a,c)} and the separating
non-link \texttt{not\ pi(a,b)}.

The active endpoint reduces \texttt{AubrunBiDefectCountBound\ m} to the
following theorem-strength inputs:

\begin{enumerate}
\def\labelenumi{\arabic{enumi}.}
\tightlist
\item
  the remaining large-gap left Biane adjacent link under
  \texttt{2\textless{}c-a} and \texttt{c-a\ !=\ 3};
\item
  the remaining large-gap right Biane adjacent link outside the checked
  \texttt{d-b=3,\ c=b+2} slice;
\item
  the side-specific residual branch collision theorems for non-fixed
  predecessor skips, together with the strict nonwrapping smaller child
  internal-subinterval collision theorem carrying the parent/child side
  data;
\item
  the Chapuy/trisection recurrence.
\end{enumerate}

Current support discipline: these live leaf groups are also rendered as
a Graphviz frontier at \texttt{build/u\_aub\_02\_frontier.svg}, tracked
in the leanblueprint declaration inventory, and checked by
\texttt{tools/build\_formalization\_artifacts.sh}. Grok review is used
only as scientific criticism of the exposition; it is not a second
formalization agent. The finite permutation scripts remain conjecture
finders only.

Previous delta: \texttt{U-AUB-349} consumes the off-cycle witness
through a smaller side-aware residual interval. The smaller
contour-repeat interval now splits: either it already gives the common
\(\pi\)/\(S\pi\) pair, or it leaves a strictly smaller
\texttt{leftTightStrictReturnIntervalResidual} with the inherited
no-\(x\)-cycle and block-off-\(x\) invariants. The current endpoint
therefore no longer asks for the whole nonwrapping off-cycle
open-internal subinterval theorem.

Most recent delta: \texttt{U-AUB-358} absorbs the strict child primitive
branches into the parent primitive residual leaf. The theorem face
therefore no longer asks separately for the child internal \texttt{S},
start \texttt{S}, end \texttt{S}, fixed predecessor, or adjacent
\texttt{T} branches. Those are all primitive residuals on a strict child
interval, hence primitive residuals on the parent interval. The
remaining smaller nonwrapping child leaf is the genuine
internal-subinterval branch.

Previous delta: \texttt{U-AUB-355} wires the strict split of the smaller
contour-repeat packet into the sharp endpoint. U-AUB-356 then replaces
raw internal \texttt{T} by adjacent \texttt{T} or genuine shorter
internal subinterval, and U-AUB-358 absorbs the primitive alternatives
as above.

Previous delta: \texttt{U-AUB-350} splits that smaller child residual by
the subinterval-preserving descent. The theorem face now asks separately
for the side-aware primitive child branch and the side-aware
internal-subinterval child branch; the raw child residual predicate is
no longer the exposed input.

Previous delta: \texttt{U-AUB-348} turns the confined off-cycle internal
witness into a strictly smaller charged contour interval. The witness
endpoints \(a+r,a+s\) form a same-\(\pi\) interval strictly inside
\((a,c)\), the interval has its own
\texttt{leftTightContourIntervalRepeat} packet, and the whole witness
block is still off the skipped \(x\)-cycle. This is support for the open
collision leaf, not a closure of that leaf.

Previous delta: \texttt{U-AUB-347} proves strict noncrossing confinement
for the off-cycle nonwrapping internal block. If \(a,c\) are the
first-return endpoints in the skipped \(x\)-cycle and the internal
witness starts at \(a+r\) off that cycle, the whole \(\pi\)-block of
\(a+r\) lies strictly inside \((a,c)\).

Previous delta: \texttt{U-AUB-329} makes the branch-normalized return
interval theorem-facing at the current sharp endpoint. The collision
theorem no longer has to consume an arbitrary minimal-frontier packet;
it can use the nonwrapping first-return, wrap-below last-return, or
wrap-above first-return interval tied to the original skip. This is a
sharper checked conditional endpoint, not a proof of the collision
theorem.

Previous delta: \texttt{U-AUB-328} completes the branch-normalized
return-interval package for nested predecessor skips. The wrap-below
branch now chooses the last same-\texttt{pi} return below \texttt{x};
because a nonfixed wrapping skip has
\texttt{x.val\ \textless{}\ (pi\ x).val}, the interval ending at
\texttt{(pi\ x).val} has a genuine gap and carries the repeated-contour
obstruction. The unified theorem packages nonwrapping first return,
wrap-below last return, and wrap-above first return. This still does not
prove the nested collision theorem. A new pure-math article source was
also created at \texttt{article/u\_aub\_02\_pure\_math\_article.tex}; it
compiles cleanly and has one Grok referee pass recorded in
\texttt{build/grok\_u\_aub\_02\_article\_referee.md}.

Previous delta: \texttt{U-AUB-327} extracts first-return contour
intervals in two branches of the nested predecessor-skip leaf. In the
nonwrapping branch, the first same-\texttt{pi} return after
\texttt{(pi\ x).val} is separated by a genuine gap unless the
already-closed direct-hit case occurs. The same first-return package is
now available for the wrap-above branch. Both packages include the
left-tight repeated-contour obstruction. They do not prove the collision
theorem.

Previous delta: \texttt{U-AUB-326} adds a tied contour extractor for the
nested predecessor-skip leaf. A non-fixed skip now produces a long
repeated contour interval together with an origin tag: either the
interval is exactly \texttt{{[}(pi\ x).val,x.val{]}}, or it is a
wrapping interval ending at \texttt{(pi\ x).val}, or it is a wrapping
interval starting at \texttt{x.val}. This prevents the later collision
proof from treating the charged interval as anonymous. It does not prove
the collision theorem.

Previous delta: \texttt{U-AUB-325} removes the direct-hit branch from
the minimal-frontier leaf. If \texttt{finRotate\ (pi\ x)} is already in
the same \texttt{pi}-cycle as \texttt{x}, the checked theorem
\texttt{permCyclicPredecessorSkip\_direct\_commonCycle} supplies the
nontrivial common \texttt{pi}/\texttt{finRotate*pi} pair. The live
frontier theorem therefore only needs to handle the nested case where
this same-cycle relation fails.

Most recent delta, right special slice: \texttt{U-AUB-320} closes the
last checked right branch, \texttt{S(b+1,d)}. If \texttt{S(b+1,d)} holds
while the target \texttt{S(c,d)} fails, uniqueness of simple incidence
paths from \texttt{S(c)} to \texttt{S(d)} forces \texttt{pi(b+1,c)}.
Combined with \texttt{pi(a,c)}, this gives \texttt{pi(a,b+1)}. The
shifted left adjacent-link input on \texttt{(a,b,b+1,d)} gives
\texttt{S(a,b)}, and then \texttt{S(a,b)}, \texttt{S(b+1,d)},
\texttt{pi(a,b+1)}, and \texttt{pi(b,d)} form the forbidden incidence
square. The active endpoint no longer exposes \texttt{hRightS\_b1\_d},
\texttt{hRightFixed\_b1}, \texttt{hRightS\_b\_c}, or
\texttt{hRightT\_b1\_c}. This is a checked closure of the right
gap-three special slice; U-AUB-02 remains open because the non-local
leaves below are still theorem-facing.

Recent delta, left end slice: \texttt{U-AUB-321} closes the fixed
\texttt{a+1} branch. In the slice \texttt{c-a=3}, \texttt{b=a+2},
fixedness of \texttt{a+1} gives the adjacent repeat \texttt{S(a+1,b)}.
If the target \texttt{S(a,b)} failed, the contour step
\texttt{a\ -\textgreater{}\ a+1}, this repeat, and the contour step
\texttt{b\ -\textgreater{}\ c} would form a simple length-four incidence
path from \texttt{S(a)} to \texttt{S(c)}. Tree uniqueness compares that
path with the direct length-two path from \texttt{pi(a,c)},
contradiction. The active endpoint no longer exposes
\texttt{hLeftEndFixed\_a1}.

Most recent delta, left end slice: \texttt{U-AUB-322} closes the
remaining \texttt{S(a+1,c)} branch. The branch copies the contour step
\texttt{a+1\ -\textgreater{}\ b} into a simple path from \texttt{S(c)}
to \texttt{S(b)}. The reversed contour step
\texttt{b\ -\textgreater{}\ c} gives a second simple path with the same
endpoints. Tree uniqueness identifies their middle right vertices, hence
forces \texttt{pi(a+1,b)}. The existing \texttt{pi(a+1,b)} absorber then
closes the target. The active endpoint no longer exposes any finite
gap-three branch.

Tooling delta: \texttt{U-AUB-323} makes the large-gap sanity profile
reproducible. The branch sanity script now checks the live large-gap
crossing hypotheses in small left-tight models; through \texttt{n=9}, it
enumerates 6,916 left-tight permutations and finds no crossing
hypotheses. This is conjecture-finding evidence only, not a substitute
for the large-gap Biane theorem.

This endpoint is a checked reduction, not a no-input closure of
U-AUB-02. Each item in the list above is still theorem-facing. Thus a
\texttt{leanok} marker for this endpoint means that the reduction from
these inputs is verified; it does not mean that the Aubrun count theorem
has been proved unconditionally. In particular, words such as
\texttt{Biane}, \texttt{frontier}, or \texttt{chapuy} inside the Lean
name are not completion markers. They name the hypotheses being
consumed.

\subsection{Mathematical route}\label{mathematical-route}

This section states the intended counting theorem, not a currently
closed checked theorem.

The corrected route is one-sided first. The intended combinatorial
theorem interprets the plus-defect \(\delta_+\) through the
unicellular-map genus: \[
  \delta_+(\pi) = 2g(\pi).
\] The expected counting strategy is then: \[
  \#\{\delta_+=0\} \le 4^{m+1},
\] and for genus \(g\), \[
  N_{m,g+1}
  \le (2(m+1))^3 N_{m,g}.
\] Iterating gives \[
  N_{m,g} \le 4^{m+1}(2(m+1))^{3g}.
\] This is the combinatorial content that should feed the bidefect
envelope. This paragraph is a proof target, not a closed checked
theorem. The current verified content is the reduction from concrete
Biane/Chapuy inputs to the bidefect envelope; the unconditional
genus-zero and Chapuy counting theorems remain open.

\subsection{Current Checked Content}\label{current-checked-content}

The proof graph contains a large amount of checked surrounding
structure.

The high plus-defect tail, finite small-window reductions, bidefect
scalar absorption, \texttt{NCPart} cardinality, gap-two Biane adjacent
links, and several minimal-contour/frontier reductions are already
verified. The latest checked endpoint sharpens the route as follows:

\begin{Shaded}
\begin{Highlighting}[]
\NormalTok{large{-}gap Biane half{-}links}
\NormalTok{+ checked gap{-}three frontier reductions}
\NormalTok{+ side{-}specific residual branch collisions}
\NormalTok{+ Chapuy recurrence}
\NormalTok{=\textgreater{} AubrunBiDefectCountBound m.}
\end{Highlighting}
\end{Shaded}

This display is a conditional implication, not a completed theorem. The
named declaration still carries the large-gap Biane inputs, the
side-specific residual return-interval collision theorems, and the
Chapuy recurrence as hypotheses.

As a conditional adapter, the endpoint narrows the previous large-gap
route in two ways. First, it removes the direct no-skip hypothesis from
the theorem face. Instead of asking directly that a noncrossing
left-tight cycle partition has no non-fixed predecessor skip, the
current endpoint asks for the sharper collision statement:

If a non-fixed predecessor skip exists, extract the side-tied long
same-\(\pi\) return interval. Split the origin side into nonwrapping
first return, wrap-below last return, and wrap-above first return; then
split each side into primitive and internal-subinterval branches. In the
nonwrapping internal-subinterval branch the endpoint-touching case and
the case where the internal witness remains in the skipped \(x\)-cycle
are checked impossible, so the theorem-facing nonwrapping child leaf is
now split into the side-aware primitive and internal-subinterval child
branches produced from the off-cycle witness; the child interval lies
strictly inside the first-return interval, has no skipped \(x\)-cycle
points inside it, and its \(\pi\)-block stays inside the parent interval
and off \(x\). The remaining branch theorems must force two distinct
points \(y,z\) such that \[
  y \sim_\pi z
  \quad\text{and}\quad
  y \sim_{\gamma\pi} z.
\] This contradicts left-tight transversality.

Second, it no longer asks the broad left large-gap Biane hypothesis to
handle any \texttt{c-a=3} slice. Since
\texttt{a\textless{}b\textless{}c} and \texttt{c-a=3} force either
\texttt{b=a+1} or \texttt{b=a+2}, both left length-three slices are now
routed through finite contour reductions.

For \texttt{c-a=3}, \texttt{b=a+1}, the left start-slice branches are
now all discharged at the active endpoint. The four former branches were
\[
  S(b,c),\quad \pi(a+2)=a+2,\quad \pi(b,a+2),
  \quad
  \bigl(\neg(a\sim_\pi a+2)\ \wedge\ S(a,a+2)\bigr).
\] The left-start \texttt{S(b,c)} among them is supplied by the right
Biane frontier. The fixed-point branch and the \texttt{pi(b,a+2)} branch
both force that already-closed left-start \texttt{S(b,c)} branch. The
last branch is closed directly by the copied length-two incidence path
described above.

For \texttt{c-a=3}, \texttt{b=a+2}, the remaining left finite branches
are \[
  S(a+1,c).
\] The earlier start-slice proof used several normalizations before the
final closure. If the raw branch \(S(a,a+2)\) also has the hidden link
\(a\sim_\pi a+2\), the checked start-gap-two right-repeat theorem gives
the target \(S(a,b)\). The adjacent branch \(S(a+2,c)\) is exactly the
adjacent repeat at \(a+2\), hence becomes the fixed-point branch
\(\pi(a+2)=a+2\). The target \(S(a,b)\) itself is impossible: \(b=a+1\),
while \(a\) already lies in a nontrivial same-\(\pi\) interval with
\(c\) and \(c-a=3\). An adjacent \((\gamma\pi)\)-repeat \(S(a,a+1)\)
would force \(a\) to be fixed by \(\pi\), contradicting \(a\sim_\pi c\)
with \(a<c\). The checked theorem
\texttt{leftTight\_bianeAdjacentLink\_left\_of\_gap\_three\_start\_adjacent\_branch\_to\_target\_iff\_false}
therefore records that branch-to-target implications in this slice are
the same as proving the corresponding branch impossible. The branch
\(\pi(a,a+1)\) is now closed: together with \(\pi(a,c)\) it gives
\(\pi(a+1,c)\), so the shifted gap-two crossing \((a+1,b,c,d)\) gives
\(S(a+1,b)\), which feeds the already exposed branch. The branch
\(\pi(a+1,b)\) is also closed: viewed on the shifted crossing
\((a,a+1,c,d)\), it is the already exposed start-adjacent \(\pi(b,a+2)\)
branch. That branch gives \(S(a,a+1)\), which is impossible for the
nontrivial same-\(\pi\) interval \(a\sim_\pi c\). The remaining
adjacent-repeat branch \(S(a+1,b)\) has also been normalized: because
\(b=a+2\), it is exactly the adjacent left repeat at \(a+1\), and
left-tightness turns it into the fixed-point branch \(\pi(a+1)=a+1\).
That fixed-point branch is now closed: fixedness gives \(S(a+1,b)\), and
a denial of \(S(a,b)\) produces a simple length-four path from \(S(a)\)
to \(S(c)\). Tree uniqueness compares it with the direct length-two path
from \(a\sim_\pi c\), contradiction.

The right branch \texttt{S(b,c)} is now closed: assuming \texttt{S(b,c)}
and denying the target \texttt{S(c,d)} gives a copied contour path from
\texttt{S(b)} to \texttt{S(d)}, and tree uniqueness against the direct
\texttt{pi(b,d)} path forces the forbidden link \texttt{pi(b,c)}. This
right special branch is no longer a theorem-facing input to the active
endpoint.

The branch \texttt{pi(b+1,c)} is also no longer independent. It combines
with the crossing link \texttt{pi(a,c)} to give \texttt{pi(a,b+1)}.
Applying the large-gap right input to the shifted crossing
\texttt{(a,b,b+1,d)} gives \texttt{S(b+1,d)}, because the shifted
crossing is not the special \texttt{c=b+2} slice.

The fixed branch \texttt{pi(b+1)=b+1} is now also internal. Fixedness
gives the adjacent repeat \texttt{S(b+1,c)}. If \texttt{S(b+1,d)}
failed, the contour step \texttt{b\ -\textgreater{}\ b+1} followed by
the copied contour step \texttt{c\ -\textgreater{}\ d} would be a simple
length-four incidence path from \texttt{S(b)} to \texttt{S(d)}. Tree
uniqueness compares that path with the direct length-two path from
\texttt{pi(b,d)}, contradiction. So the fixed branch feeds the same
\texttt{S(b+1,d)} implication.

The last right branch \texttt{S(b+1,d)} is now closed. If the target
\texttt{S(c,d)} were false, path uniqueness would force
\texttt{pi(b+1,c)}. Then \texttt{pi(a,c)} gives \texttt{pi(a,b+1)}, the
shifted left link gives \texttt{S(a,b)}, and the four incidences
\texttt{pi(a,b+1)}, \texttt{pi(b,d)}, \texttt{S(a,b)}, \texttt{S(b+1,d)}
form a forbidden square.

There are no remaining right finite branches in the checked
\texttt{d-b=3}, \texttt{c=b+2} slice. The right branch \(\pi(b,c)\) is
now closed: together with the crossing link \(\pi(a,c)\), it would imply
\(\pi(a,b)\), contradicting the separating non-link. The adjacent-repeat
branch \(S(b+1,c)\) has also been normalized: because \(c=b+2\), it is
exactly the adjacent left repeat at \(b+1\), and left-tightness turns it
into the fixed-point branch \(\pi(b+1)=b+1\).

These branch lists are endpoint-local snapshots. The same public
endpoint still separately requires the large-gap left Biane input, the
large-gap right Biane input, the side-specific residual collision
theorems, and the Chapuy recurrence.

\subsection{What remains open}\label{what-remains-open}

U-AUB-02 is not done. The remaining theorem-strength leaves are:

\begin{enumerate}
\def\labelenumi{\arabic{enumi}.}
\tightlist
\item
  prove the remaining large-gap left Biane adjacent-link cases;
\item
  prove the remaining large-gap right Biane adjacent-link cases;
\item
  prove the remaining side-specific residual collision theorems, now
  with branch-normalized return data available in all three nested
  branches and the nonwrapping child residual carrying the inherited
  no-\(x\)-cycle and block-off-\(x\) data; only the genuine smaller
  child internal-subinterval branch remains theorem-facing after child
  primitive absorption;
\item
  prove the Chapuy/trisection recurrence.
\end{enumerate}

These branches are not mere notation. They are finite Biane contour
cases that likely need incidence-tree uniqueness or a descent lemma, not
a naive cycle-transitivity shortcut. The finite branch scripts are
conjecture finders only; their clean output is not counted as proof.

\subsection{Checked reduction
evidence}\label{checked-reduction-evidence}

The new endpoint is in:

\begin{Shaded}
\begin{Highlighting}[]
\NormalTok{PptFactorization/AppendixBAubrunGraduate.lean}
\end{Highlighting}
\end{Shaded}

The active checked declaration is:

\begin{Shaded}
\begin{Highlighting}[]
\NormalTok{PptFactorization.AppendixB.AubrunBiDefectCountBound\_of\_bianeAdjacentLinkHalves\_left\_gap\_three\_start\_impossible\_frontiers\_nestedReturnContourIntervalCommonCycle\_and\_chapuy}
\end{Highlighting}
\end{Shaded}

The previous reduction layer is:

\begin{Shaded}
\begin{Highlighting}[]
\NormalTok{PptFactorization.AppendixB.AubrunBiDefectCountBound\_of\_bianeAdjacentLinkHalves\_left\_gap\_three\_frontiers\_skipMinimalFrontierCommonCycle\_and\_chapuy}
\end{Highlighting}
\end{Shaded}

Verification performed:

\begin{Shaded}
\begin{Highlighting}[]
\NormalTok{python3 scripts/lean\_quick\_check.py PptFactorization.AppendixBAubrunGraduate {-}{-}timeout 300}
\NormalTok{lake build PptFactorization.AppendixBAubrunGraduate}
\NormalTok{lake build PptFactorization.AristotleTargets.UpperFromLowerColumnConcreteChoices}
\NormalTok{python3 scripts/aubrun\_branch\_sanity.py {-}{-}max{-}n 9 {-}{-}all}
\NormalTok{Pandoc rebuild of exposition/u\_aub\_02\_formalization\_note.pdf}
\NormalTok{Blueprint print{-}PDF rebuild and declaration inventory check}
\NormalTok{Grok adversarial review}
\NormalTok{\#print axioms ...left\_gap\_three\_start\_impossible\_frontiers\_nestedReturnContourIntervalCommonCycle\_and\_chapuy}
\NormalTok{\#print axioms ...left\_of\_gap\_three\_end\_adjacent\_fixed\_a1}
\end{Highlighting}
\end{Shaded}

The axiom audit reports only:

\begin{Shaded}
\begin{Highlighting}[]
\NormalTok{[propext, Classical.choice, Quot.sound]}
\end{Highlighting}
\end{Shaded}

This audit is evidence that the current reduction layer has no
project-specific axioms. It is not evidence for the open leaves
themselves. In particular, the large-gap Biane links, the nested
non-direct return-interval collision theorem, and the Chapuy recurrence
still need their own checked declarations and axiom audits.

\subsection{Tooling status}\label{tooling-status}

Automatic prose narration is not used as proof evidence. The current
note is hand-curated from the checked Lean declarations, and any future
machine-written prose must be audited against the exact hypothesis list
of the active endpoint.

Pandoc is installed and converts this note to both LaTeX and PDF using
the repository TeX setup. This PDF has been refreshed from the current
Markdown. Leanblueprint was initialized for the project, with a
declaration inventory at \texttt{blueprint/lean\_decls}; blueprint
refreshes are treated as a separate verification step and are not proof
evidence by themselves. Full doc-gen integration was deliberately
deferred when it attempted to download a 518 MB Lean release-candidate
toolchain, which would slow the current proof workflow.

A finite branch sanity script is now available at
\texttt{scripts/aubrun\_branch\_sanity.py}. It enumerates small
left-tight permutations and checks active branch implications for
counterexamples before they are attempted in Lean. This is only a
guardrail: the current script reports no active finite gap-three
branches, then records the large-gap profile through \texttt{n=9}.

Grok was used for advisory proof-architecture and exposition review. Its
commentary remains advisory; Lean declarations, builds, axiom audits,
Pandoc artifacts, and blueprint checks are the evidence.

\end{document}
